\section{Continuous references, held-action moments, and output-error identities}\label{r11:proof}
This appendix proves Proposition~\ref{r11:decomposition} and specifies the complete arithmetic used in Section~\ref{r11:accuracy}. It is independent of the stopped-economy dual. All quantities here concern the linear-quadratic objective \eqref{r11:lq}; no claim about a constrained nonquadratic optimum is obtained by replacing its payoff with this one.

\subsection{Continuous verification and the closed reference}
We give a scalar proof; summation and division by $d$ yield the stated result. The finite-horizon Riccati equation with $q\geq0$, $r,w>0$ has a nonnegative bounded solution. This can also be checked directly from the displayed reference formula: it is the fractional transformation generated by
\[
 \frac{d}{ds}\begin{pmatrix}A(s)\\B(s)\end{pmatrix}
 =\begin{pmatrix}-a&c\\q&a\end{pmatrix}\begin{pmatrix}A(s)\\B(s)\end{pmatrix},
 \qquad (A(0),B(0))=(1,w),\quad c=b^2/r.
\]
The matrix square is $(a^2+cq)I$, giving the hyperbolic formula in the text. With $P(s)=B(s)/A(s)$, differentiation gives $P'=q+2aP-cP^2$. Nonnegative initial data stay nonnegative, and the negative quadratic term prevents a finite-time upward explosion. Moreover $A'=(-a+cP)A$ and $A(0)=1$, so $A$ stays positive and
\[
 \int_0^1P(s)ds=\frac{\log A(1)+a}{c}.
\]
At $\lambda=0$ in the specified cases, $a=q=0$, $A(s)=1+cws$, and $B(s)=w$. This branch is evaluated by its exact limit, not by division by a small interval containing zero.

Set $v(t,x)=p(t)x^2+\sigma^2\int_t^1p(s)ds$, where $p(t)=P(1-t)$. Direct substitution gives
\[
 v_t+(ax+bu)v_x+\tfrac12\sigma^2v_{xx}+qx^2+ru^2
 =r\left(u+\frac{bp(t)}r x\right)^2.
\]
For a square-integrable adapted control on a finite horizon, the linear state equation has finite integrated and terminal second moments. Localizing the It\^o integral and passing to the limit in the nonnegative quadratic terms gives the verification identity \eqref{r11:continuousdefect}; equivalently one can first prove it for bounded stopped controls and use the explicit linear state solution and $L^2$ approximation. The feedback $u=-bp(t)X_t/r$ is admissible, has finite second moments, and makes the square vanish. It therefore attains the value in \eqref{r11:reference}. This proves both directions of optimality rather than using a feasible Riccati policy as a lower bound without justification.

\subsection{Exact deployment coefficients}
For one interval of length $h$, put
\[
 E=e^{ah},\quad I=\int_0^h e^{as}ds,\quad A_2=\int_0^h e^{2as}ds,
 \quad B_2=\int_0^h e^{as}\!\left(\int_0^s e^{av}dv\right)ds,
\]
\[
 C_2=\int_0^h\left(\int_0^s e^{av}dv\right)^2ds,\qquad
 D_2=\int_0^h\int_0^s e^{2av}dv\,ds.
\]
For $a\ne0$, these quantities are
\begin{align*}
 I&=(e^{ah}-1)/a,& A_2&=(e^{2ah}-1)/(2a),\\
 B_2&=(A_2-I)/a,& C_2&=(A_2-2I+h)/a^2,&D_2&=(A_2-h)/(2a).
\end{align*}
For $a=0$, they are $I=A_2=h$, $B_2=D_2=h^2/2$, and $C_2=h^3/3$. Sixty-decimal interval evaluation, with a separate zero-drift branch, encloses cancellation in these expressions. All actual nonzero drift inputs in the finite suite are point intervals bounded away from zero; a division by an interval containing zero is rejected rather than repaired numerically.

Conditional on $X_{ih}=x$ and an action $u$ held through the interval,
\[
 \E[X_{(i+1)h}^2\mid x,u]=(Ex+bIu)^2+\sigma^2A_2,
\]
\[
 \E\!\left[\int_{ih}^{(i+1)h}(qX_t^2+ru^2)dt\,\middle|\,x,u\right]
 =qA_2x^2+2qbB_2xu+(qb^2C_2+rh)u^2+q\sigma^2D_2.
\]
The formulas retain the correlation between the held action and the state at the beginning of the slab. They do not approximate the action by a continuous feedback evaluated along the evolving state.

If the remaining value after this slab is $P_{i+1}x^2+z_{i+1}$, define
\begin{align*}
 D_i&=qb^2C_2+rh+P_{i+1}b^2I^2>0,\\
 L_i&=qbB_2+P_{i+1}EbI,\qquad K_i^*=L_i/D_i,\\
 P_i&=qA_2+P_{i+1}E^2-L_i^2/D_i,\\
 z_i&=z_{i+1}+q\sigma^2D_2+P_{i+1}\sigma^2A_2,
 \qquad P_N=w,\ z_N=0.
\end{align*}
The Bellman minimization is exact completion of the square. Nonnegativity of costs implies $P_i\geq0$, so the strict inequality for $D_i$ follows from $rh>0$. It applies to the entire unbounded action line, not merely to a stationary point. These $N$ backward steps are closed exact evaluation and improvement operations for quadratic values and linear held-action outputs.

For any adapted held-action rule, its conditional Bellman excess at stage $i$ is $D_i(u_i+K_i^*X_{ih})^2$. Telescoping the conditional expectations gives
\[
 J(u)-J_h^*=\frac1d\sum_{i,j}D_{ij}\E[(u_i^j+K_{ij}^*X^j_{ih})^2].
\]
Combining this equality with the continuous verification result proves \eqref{r11:twodefect}. In particular, for an implemented linear actor $u_i^j=-K_{ij}X^j_{ih}$, the optimization error is weighted by the \emph{deployed actor's} second moments, not by moments under an ideal policy. If $|K_{ij}-K_{ij}^*|\leq\delta_{ij}$ and these moments are at most $M_{ij}$, the explicit allowance is $d^{-1}\sum D_{ij}\delta_{ij}^2M_{ij}$. This is an output-level certificate for the implemented optimizer; it is not a statement that its weight-space objective is convex.

\subsection{Independent interval policy evaluation}
The verifier does not reuse the backward value as the value of a stored floating-point policy. Instead, starting with $m_0=x_0^2$, it propagates
\begin{align*}
 m_{i+1}&=(E-bIK_i)^2m_i+\sigma^2A_2,\\
 \Delta J_i&=\{q(A_2-2bK_iB_2+b^2K_i^2C_2)+rhK_i^2\}m_i+q\sigma^2D_2,\\
 J(K)&=\sum_i\Delta J_i+w m_N.
\end{align*}
Each coefficient and stored gain is imported as the exact binary64 value that defines the deployed model and actor. The arithmetic is outward at 60 decimal digits. Consequently the policy-cost interval and the closed-reference interval enclose two independently computed original-model quantities. Subtraction in interval arithmetic yields an absolute-regret interval; a negative upper endpoint is an error and is not clipped to zero. Conversion to machine-readable binary64 endpoints expands the lower endpoint toward $-\infty$ and the upper endpoint toward $+\infty$ by one unit. The output files hash the little-endian binary64 gain array, binding every certificate to its policy.

The differentiable Adam objective is a separate torch implementation of the moment equations in binary64. Its agreement with the interval evaluator is checked to a small conversion tolerance, but that agreement is only a software cross-check. The interval equations and continuous-reference proof, not the tolerance in that test, establish the certificate. Additional tests check the closed formula against a separately integrated Riccati equation, test the scalar tracking formula, reject nonfinite gains, and verify that a deliberately perturbed actor changes both its policy hash and cost. Timings and arithmetic intervals are kept separate: timing fluctuations cannot change whether a stored policy is certified at the prescribed tolerance.
